\uuid{K5g9}
\titre{Matrice d'une application linéaire/Espaces supplémentaires/Polynômes}
\auteur{hanani ; mzari}
\organisation{exo7}
\licence{CC BY-NC-SA 4.0}
\datecreate{2020-01-20}
\niveau{L1}
\module{Algèbre}
\chapitre{Matrice}
\sousChapitre{Matrice et application linéaire}
\difficulte{3}
\theme{Matrice d'une application linéaire, Espaces supplémentaires, Polynômes}
\exo7qcmid{lille-2.matrice-application.015}

\contenu{
\begin{question}
\qclassification{108.03}
\qtags{motcle=Matrice d'une application linéaire/Espaces supplémentaires/Polynômes}
On considère $\Rr_3[X]$, l'ensemble des polynômes à coefficients réels de degré $\le 3$, muni de sa base canonique notée ${\cal {B}}$ et $f$ l'application linéaire définie par :
$$\begin{array}{rccc}f:&\Rr_3[X] &\to& \Rr_3[X]\\
& P&\to &R, \end{array}$$
où $R$ est le reste de la division euclidienne de $P$ par $(X-1)^2$. Soit ${\cal {B'}} =\{P_1,P_2,P_3,P_4\}$, où $P_1=1$, $P_2=1-X$, $P_3=(1-X)^2$ et $P_4=X(1-X)^2$. Quelles sont les assertions vraies ?
\begin{answers}
\good{La matrice de $f$ dans la base $ {\cal {B}}$ est : $\left(\begin{array}{rccc}
1&0&-1&-2\\
0&1&2&3\\
0&0&0&0\\
0&0&0&0\\
\end{array}\right)$.}
\good{${\cal {B'}} $  est une base de $\Rr_3[X]$}.
\bad{La matrice de $f$ dans la base $ {\cal {B'}}$ est : $\left(\begin{array}{rccc}
0&0&1&0\\
0&0&0&1\\
0&0&0&0\\
0&0&0&0\\
\end{array}\right)$.}
\good{$\ker f$ et $\Im f$ sont supplémentaires dans  $\Rr_3[X]$.}
\end{answers}
\begin{explanations} La matrice de $f$  d'une base $ {\cal {B}}=(u_j)$ dans une autre base $ {\cal {B}}'=(v_i)$ est la matrice  dont la $j$ième colonne est constituée des coordonnées de  $f(u_j)$ dans la base $ {\cal {B}}'$.
\vskip0mm
La matrice de $f$ dans la base ${\cal {B}}$ est : $\left(\begin{array}{rccc}
1&0&-1&-2\\
0&1&2&3\\
0&0&0&0\\
0&0&0&0\\
\end{array}\right)$.
\vskip0mm
La matrice de $f$ de la base $ {\cal {B'}}$ est : $
\left(\begin{array}{rccc}
1&0&0&0\\
0&1&0&0\\
0&0&0&0\\
0&0&0&0\\
\end{array}\right).$
\vskip0mm
De cette matrice, on déduit que $\{P_1,P_2\}$ est une base de $\Im f$ et $\{P_3,P_4\}$ est une base de $\ker f$. Comme $\{P_1,P_2,P_3,P_4\}$ est une base de $\Rr_3[X]$, $\Im f$ et $\ker f$ sont  supplémentaires dans  $\Rr_3[X]$.
\end{explanations}
\end{question}
}
